\chapter{Conclusion}
\label{ch: conclusion}


Flooding is set to remain a lasting challenge in France. A recent report by Réseau Action Climat (RAC) and Ademe \cite{RAC2024} provides a detailed region-by-region assessment of France’s changing climate landscape, highlighting diverse climatic risks and vulnerabilities. Compared to the IPCC’s global outlook, these assessments are more local in scale, and thus directly relevant for adaptation: they fill knowledge gaps where policy must operate. France’s \emph{Plan national d’adaptation au changement climatique} (PNACC), whose Measure 3 explicitly targets flooding, is one such framework, but its implementation has often been uneven and interrupted—underscoring the difficulty of sustaining long-term climate governance.

At the European level, recent debates—such as those raised in the Draghi Report—stress that climate policy cannot be separated from industrial policy. A green industrial strategy is increasingly necessary to overcome adaptation and mitigation barriers. In this regard, Earth observation systems such as Copernicus provide comprehensive, real-time data for environmental and climate monitoring, while new investments in green infrastructure and innovation are central to Europe’s adaptation capacity.

It is important to recognize that the very notion of a “natural disaster” is not neutral. As the vulnerability paradigm emphasized since the 1970s, disasters are not solely the result of extreme environmental events but also of socio-economic and political conditions that make certain groups more susceptible to harm \citep{nations2010natural}. The IPCC’s Sixth Assessment Report similarly stresses that changes in the frequency or magnitude of hazards do not directly translate into changes in risk if proactive management—such as protection or managed retreat—is undertaken \citep{Reisinger2020Concept}. In this sense, what appears as a sudden disaster is in fact the culmination of long-term exposure, behavior, and institutional choices. Climate change exacerbates these dynamics by introducing non-stationarity into natural systems, making past data a less reliable guide for future adaptation.

The urgency of this issue is underscored by recent events. The European State of the Climate 2024 reports that nearly one-third of Europe’s river network flooded last year \citep{copernicus2024esotc}. In France, October 2024 floods in the Auvergne–Rhône–Alpes region brought 600 mm of rain in just 48 hours—roughly Paris’s annual total—triggering widespread evacuations, crippling infrastructure, and leading the Ministry for Ecological Transition to relaunch the PNACC. Meteo France described the episode as the most intense rainfall over two days since the early 20th century. Similar catastrophic floods in Spain, Germany, and the Netherlands further illustrate the rising risks posed by climate extremes across Europe.

In light of the empirical evidence presented in this paper—showing that floods generate substantial microeconomic disruptions but muted aggregate effects—the conclusion is clear: institutional design matters. Insurance, fiscal buffers, spatial planning, and public information can prevent localized shocks from cascading into systemic crises. Yet climate change is intensifying the scale and frequency of hazards, and adaptation policies must be calibrated to fat-tailed risks, not historical averages. The challenge for France and Europe is to ensure that climate adaptation is not reactive but anticipatory, linking climate policy, industrial policy, and social protection. Doing so will not only reduce aggregate losses but also, crucially, limit the unequal distribution of harm across communities, workers, and firms.



% Flooding in France is here to remain, \href{https://reseauactionclimat.org/la-france-face-au-changement-climatique-toutes-les-regions-impactees/}{report by Réseau Action Climat (RAC)} and Ademe provides a detailed region-by-region assessment of France's changing climate landscape, highlighting diverse climatic risks and vulnerabilities. More local than IPCCC, so this is also a magin of adaptation: filling knowledge gaps. Also, the Plan national d'adaptation au changement climatique (PNACC) Measure 3 directly concerns flooding and it says.. 




% Draghi rport:, and industrial strategy for Europe to overcome these barriers; applying a climate policy without an industrial policy, indutriall policis on the rise  

% In Earth Observation, Copernicus offers the most comprehensive
% data worldwide, including for environmental and climate change monitoring, disaster management and security.


% Economics of catastrophes: Global Priorities Institute \\

% Science Po Syllabus \\


% We should not take for granted that it is a disasdter, but ask: What made it so? In Switzerlannd, Blatten was not a disaster because all were evvacuated!\\

% Ethimology of ``disaster" pag. 54 of pdf Disaster Studies \\

% A green innovation/industrial policy for Europe (Reinhilde Veugelers) \\




% \textbf{The notion of natural disaster:} Because both reactive and proactive actions can influence the ultimate impact of natural disasters, the notion itself of “natural disaster” is an analytical conceit that needs interpretation. As a rupture of the socio-economical order, a natural disaster is triggered by environmental forces but is ultimately an historical process rooted in the path-dependent ways economic and governance institutions change. This recognition led to the emergence in the 1970s of the \textit{vulnerability} paradigm in reaction to the then dominant \textit{hazard} paradigm, shifting the focus from the hazard event itself to the underlying conditions that make certain groups more susceptible to harm, highlighting how natural disasters are the cumulative implications of many earlier decisions, some taken individually, others collectively \citep{nations2010natural}. Hence the idea of “natural Hazards, unnatural disasters. In the IPCC's Sixth Assessment Report the notion of Risk \citep{Reisinger2020Concept} is pointed out that changes in the frequency and/or magnitude of climatte hazards do not map directly changes in actual risk if socio-economic changes like proactive risk management decisions such as protection or managed retreat are met. These is an important point to keep in mind when considering the distributional impacts of clilmate adaptation policies, as well as the estimation of damages from disasters. To conclude, what appears as a sudden “disaster” is in fact the result of an historical process deeply embedded in pre-existing patterns of exposure and behavior. As a consequence, structural exposure to natural disasters is not random, rather is rooted in long-term social and spatial processes. Climate change further exacerbates these issues as it introduces non-stationary in the distributions of many natural systems and makes climate patterns no longer stable or predictable based on past trends, a point we will return to later. \\



% \textbf{Recent floods in Europe}
% The \emph{European State of the Climate 2024} reports that almost one-third of the continent’s river network flooded last year \citep{copernicus2024esotc}, and recent deadly floods across Europe are a reminder of the growing threat of climate-induced extreme weather events. 
% Recent floods in France in October 2024 in the Auvergne-Rhône-Alpes region, delivering up to 600 mm of rain in 48 hours in the Ardèche departmeent only roughly Paris’s annual total triggering 900 evacuations and crippling rail and power infrastructure emmploying 3 000 sapeurs-pompiers and the ministry for Ecological transition called for the Relaunch of the National Climate Change Adaptation Plan (PNACC) previously stopped. The most intense ever recorded over two days since the beginning of the 20th, said Meteo France! \href{https://www.franceinfo.fr/replay-radio/le-grand-temoin/intemperies-des-moyens-massifs-toujours-deployes-selon-le-ministre-delegue-a-la-securite-du-quotidien_6816923.html}{France Info}. In neighouringh european countriesr such as Spain, Germany, and the Netherlands recent destructive floods are a reminder of the growing threat of climate change-induced extreme weather events. 

% Banchmark across countries and shared learning is very important: netherlands has an adanced flood anagmetnn sstemms due to histotry, and U we can learn: they do sponge parks and rain tunnel, sopraeleveate houses (in france evelstion of flood walls, works performed on th ebed of rivers, and the creation of reservoir lakes..).
